%macbook9T5 
\setcounter{ex}{0}
\setcounter{paragraph}{0}
\PhanI
\loai{Biên độ cách nút 1 đoạn x}
\begin{ex}
Quan sát trên một sợi dây thấy có sóng dừng với biên độ của bụng sóng là a. Tại điểm trên sợi dây cách bụng sóng một phần tư bước sóng có biên độ dao động bằng
\choice
{$\dfrac{a}{2}$}
{\True $0$}
{$\dfrac{a}{4}$}
{$a$}
\loigiai{}
\end{ex}
\begin{ex}%[MB9T5-Add1][Loai1]
Trên một sợi dây đang có sóng dừng, biên độ dao động tại các bụng là $a$.  Một điểm $M$ dao động với biên độ $0,6a$ và vị trí cân bằng của $M$ cách nút gần nhất $d=5\ cm$. Bước sóng của sóng trên dây bằng
\choice
{40,0 cm}
{46,2 cm}
{\True 49,0 cm}
{55,5 cm}
\loigiai{
\begin{itemize}
\item Biên độ tại $M$: $A_M=a\left|\sin\dfrac{2\pi d}{\lambda}\right|=0,6a$  
\item $\Rightarrow\;\sin\dfrac{2\pi d}{\lambda}=0,6\;\Longrightarrow\;\dfrac{2\pi d}{\lambda}=\arcsin0,6$  
(chọn giá trị nhỏ nhất thuộc $(0;\pi/2)$).  
\item Suy ra $\lambda=\dfrac{2\pi d}{\arcsin0,6}\approx\dfrac{2\pi\cdot5}{0,6435}\approx49,0\ cm$.
\end{itemize}}
\end{ex}
\begin{ex}%[3P2K9-2][Loai1]
Một sợi dây AB có chiều dài 1 m căng ngang, đầu A cố định, đầu B gắn với một nhánh của âm thoa. Trên dây AB có một sóng dừng ổn định với 4 bụng sóng, biên độ bụng sóng là 2 cm, B được coi là nút sóng. Điểm trên dây có vị trí cân bằng cách A một đoạn $\dfrac{13}{24}$ m dao động với biên độ là
\choice
{\True $1\ cm$}
{$2\ cm$}
{$\sqrt{2}\ cm$}
{$\sqrt{3}\ cm$}
\loigiai{
\begin{itemize}
\item Dây cố định hai đầu có 4 bụng nên $\ell =4\dfrac{\lambda }{2}=1\ m\Rightarrow \lambda =0,5\ m$.
\item Điểm trên dây có VTCB cách nút A: $x = \dfrac{13}{24}m$ suy ra biên độ ${A_M}=A_B. \left| \sin \dfrac{2\pi x}{\lambda } \right|=2\left| \sin \dfrac{2\pi \dfrac{13}{24}}{0. 5} \right|=1\ cm$.
\end{itemize}
}
\end{ex} \haidong{20}
\begin{ex}%[MB9T5-Add2][Loai1]
Một sợi dây căng ngang với đầu $A$ cố định (nút) đang có sóng dừng.  
Hai điểm $P$ và $Q$ có vị trí cân bằng cách $A$ lần lượt $x_P=12\ cm$ và $x_Q=22\ cm$.  
Biết sóng trên dây có bước sóng $\lambda=40\ cm$.  
Tỉ số giữa biên độ dao động của $P$ và $Q$ là
\choice
{$\dfrac{3}{2}$}
{$2$}
{\True $3,08$}
{$0,50$}
\loigiai{
\begin{itemize}
\item Đầu $A$ là nút $\Rightarrow$ biên độ tại điểm cách $A$ đoạn $x$:  
$A(x)=2a\left|\sin\dfrac{2\pi x}{\lambda}\right|$.
\item $A_P=\left|\sin\dfrac{2\pi\cdot12}{40}\right|=\left|\sin0,6\pi\right|\approx0,9511$.
\item $A_Q=\left|\sin\dfrac{2\pi\cdot22}{40}\right|=\left|\sin1,1\pi\right|\approx0,3090$.
\item $\displaystyle\dfrac{A_P}{A_Q}\approx\dfrac{0,9511}{0,3090}\approx3,08$.
\end{itemize}}
\end{ex}
\begin{ex}%[MB9T5-Add3][Loai1]
Trên dây $AB$ có sóng dừng với bước sóng $\lambda$. Tại một bụng $B$ biên độ dao động là $A_b = 5\ cm$. Điểm $M$ gần $B$ nhất có biên độ $3\ cm$; điểm $N$ gần $M$ nhất (về phía xa $B$) có biên độ $2\ cm$. Khoảng cách $MN$ bằng
\choice
{$\dfrac{\lambda}{16}$}
{$\dfrac{\lambda}{12}$}
{\True $\dfrac{\lambda}{20}$}
{$\dfrac{\lambda}{8}$}
\loigiai{
\begin{itemize}
\item Lấy $B$ làm gốc tọa độ ($x=0$) với $0<x<\dfrac{\lambda}{4}$, biên độ tại $x$:  
$A(x)=A_b\left|\cos\dfrac{2\pi x}{\lambda}\right|$.
\item Với $A_M=3\;$cm: $\displaystyle\cos\dfrac{2\pi x_M}{\lambda}=\dfrac{3}{5}\;\Longrightarrow\;x_M=\dfrac{\lambda}{2\pi}\arccos\dfrac{3}{5}$.
\item Với $A_N=2\;$cm: $\displaystyle\cos\dfrac{2\pi x_N}{\lambda}=\dfrac{2}{5}\;\Longrightarrow\;x_N=\dfrac{\lambda}{2\pi}\arccos\dfrac{2}{5}$.
\item Khoảng cách $MN=x_N-x_M=\dfrac{\lambda}{2\pi}\!\left(\arccos\dfrac{2}{5}-\arccos\dfrac{3}{5}\right)\!=\dfrac{\lambda}{20}$.
\end{itemize}}
\end{ex}
\loai{Tính khoảng cách}
\begin{ex}%[3P2K9-2][Loai1]
Thí nghiệm sóng dừng trên một sợi dây có hai đầu cố định và chiều dài 36 cm, người ta thấy có 6 điểm trên dây dao động với biên độ cực đại. Khoảng thời gian ngắn nhất giữa hai lần dây duỗi thẳng là 0,25 s. Khoảng cách từ bụng sóng đến điểm gần nó nhất có biên độ bằng nửa biên độ của bụng sóng là
\choice
{4 cm}
{\True 2 cm}
{3 cm}
{1 cm}
\loigiai{
\begin{itemize}
\item $\ell = \dfrac{{6\lambda }}{2} \Rightarrow \lambda = 12cm;\dfrac{T}{2} = 0,25 \Rightarrow T = 0,5\ s$.
\item ${a_M} = {a_b}cos\dfrac{{2\pi d}}{\lambda } = \dfrac{{{a_b}}}{2} \Rightarrow cos\dfrac{{2\pi d}}{\lambda } = \dfrac{1}{2} \Rightarrow \dfrac{{2\pi d}}{\lambda } = \dfrac{\pi }{3} \Rightarrow d = \dfrac{\lambda }{6} = 2\ cm$.
\end{itemize}
}
\end{ex} \haidong{20}
\loai{Trạng thái giữa hai điểm trên cùng 1 bó sóng}
\begin{ex}%[3P2K9-2][Loai1]
Một sợi dây đàn hồi căng ngang, đang có sóng dừng ổn định. Khoảng thời gian giữa hai lần liên tiếp sợi dây duỗi thẳng là 0,1 s, tốc độ truyền sóng trên dây là 3 m/s. Khoảng cách giữa hai điểm gần nhau nhất trên sợi dây dao động cùng pha và có biên độ dao động bằng một nửa biên độ của bụng sóng là
\choice
{\True 20 cm}
{30 cm}
{10 cm}
{8 cm}
\loigiai{
\begin{itemize}
\item $T = 2. 0,1 = 0,2$ s.
\item Bước sóng : $\lambda = v. T = 0,6\ m = 60$ cm.
\item Các điểm trong cùng một bó sóng dao động cùng pha.
\item Phương trình sóng dừng tại M cách nút N một khoảng d: $u=2a\cos (\dfrac{2\pi d}{\lambda }+\dfrac{\pi }{2})\cos (\omega t-\dfrac{\pi }{2})$\\ 
$A_M = 2a \cos( \dfrac{2\pi d}{\lambda } + \dfrac{\pi }{2} ) = a \Rightarrow \cos( \dfrac{2\pi d}{\lambda } + \dfrac{\pi }{2} ) = \dfrac{1}{2}$\\ 
$\Rightarrow \dfrac{2\pi d}{\lambda } + \dfrac{\pi }{2} = \pm \dfrac{\pi }{3} + k\pi \Rightarrow d = (\pm \dfrac{1}{6} - \dfrac{1}{4} + \dfrac{k}{2} )\lambda$\\ 
$\Rightarrow d_1 = (- \dfrac{1}{6} - \dfrac{1}{4} + \dfrac{k}{2} )\pi \Rightarrow d_{1\min} = (- \dfrac{1}{6} - \dfrac{1}{4} + \dfrac{1}{2} )\lambda \Rightarrow d_{1\min} = \dfrac{\lambda }{12}$\\ 
$\Rightarrow d_2 = ( \dfrac{1}{6} - \dfrac{1}{4} + \dfrac{k}{2} )\lambda \Rightarrow d_{2\min} = ( \dfrac{1}{6} - \dfrac{1}{4} + \dfrac{1}{2} )\lambda \Rightarrow d_{2\min} = \dfrac{5\lambda }{12}$\\ 
$MM’ = d_{2\min} - d_{1\min} = \dfrac{5\lambda }{12} - \dfrac{\lambda }{12} = \dfrac{\lambda }{3} = 20$ cm .
\end{itemize}
}
\end{ex} \haidong{20}
\loai{Hai điểm có cùng biên độ}
\begin{ex}%[3P2K9-2][Loai1]
Một sợi dây có sóng dừng hai đầu cố định với tần số 5 Hz. Biên độ dao động của điểm bụng là 2 cm. Khoảng cách gần nhất giữa hai điểm trên hai bó sóng cạnh nhau có cùng biên độ 1 cm là 2 cm. Tốc độ truyền sóng là
\choice
{1,2 m/s}
{0,8 m/s}
{\True 0,6 m/s}
{0,40 m/s}
\loigiai{
\begin{itemize}
\item Ta có ${A_M}=A_b. \left| \sin \dfrac{2\pi x}{\lambda } \right|=A_b\left| \sin \dfrac{2\pi fx}{v} \right| \left(0 <\dfrac{2\pi x}{\lambda} \le \dfrac{\pi }{2}\right)$.
\item Vì x là tọa độ dương của điểm gần nhất thì có $0 < x \le \dfrac{\lambda }{4}$ ) hay $1=2\left| \sin \dfrac{2\pi . 5. 0,01}{v} \right|\Rightarrow \sin \dfrac{0,1\pi }{v}=\pm \dfrac{1}{2}\Rightarrow v=0,6\ m/s$.
\end{itemize}
}
\end{ex} \haidong{20}
\begin{ex}
Tạo sóng dừng trên sợi dây $PQ = 55\ cm$, $P$ là nút, $Q$ là bụng, $f = 100\ Hz$. Có 5 nút khác ngoài $P$. Độ rộng bụng sóng là $10\ mm$. Gia tốc cực đại tại điểm $M$ cách $Q$ $2,5\ cm$ bằng
\choice
{$4,2\ m/s^2$}
{$3,2\ m/s^2$}
{$1,1\ m/s^2$}
{\True $2,2\ m/s^2$}
\loigiai{
Có 6 nút sóng $\Rightarrow PQ = \dfrac{11}{2} \lambda = 55 \Rightarrow \lambda = 20\ cm$\\
$A_M = A_b |\cos \dfrac{2\pi . MQ}{\lambda}| = 5 . |\cos \dfrac{2\pi . 2,5}{20}| = 2,5\sqrt{2}\ mm$\\
$a_{M\max} = \omega^2 A_M = (2\pi f)^2 A_M \approx 2,2\ m/s^2$
}
\end{ex}
\begin{ex}
Thí nghiệm sóng dừng trên sợi dây với tần số $30\ Hz$, có 2 bó sóng, chiều dài mỗi bó $10,2\ cm$, bề rộng mỗi bó $5,6\ cm$. Tỉ số giữa tốc độ truyền sóng và tốc độ cực đại điểm bụng là
\choice
{\True 1,2}
{2,1}
{3,2}
{1,6}
\loigiai{
$\lambda = 2 . 10,2 = 20,4\ cm$, $A_b = \dfrac{5,6}{2} = 2,8\ cm$\\
$\dfrac{v}{v_{b\max}} = \dfrac{\lambda f}{2\pi f A_b} = \dfrac{\lambda}{2\pi A_b} = \dfrac{20,4}{2\pi . 2,8} \approx 1,2$
}
\end{ex}
\begin{ex}%[3P2K9-2][Loai1]
Trên một sợi dây có sóng dừng với biên độ điểm bụng là 5 cm. Giữa hai điểm M và N trên dây có cùng biên độ dao động 2,5 cm, cách nhau 20 cm các điểm luôn dao động với biên độ nhỏ hơn 2,5 cm. Bước sóng trên dây là
\choice
{\True 120 cm}
{80 cm}
{60 cm}
{40 cm}
\loigiai{
\begin{itemize}
\item M và N cách đều nút 1 đoạn : $d = 10$ cm, ta có : $a_M = 2a\sin(2\pi \dfrac{d}{\lambda } )\ (2a = 5\ cm) \Rightarrow \sin(2\pi \dfrac{d}{\lambda }) = \dfrac{1}{2} \Rightarrow 2\pi \dfrac{d}{\lambda }= \dfrac{\pi}{6} \Rightarrow d = \dfrac{\lambda}{12} \Rightarrow \lambda = 120$ cm.
\end{itemize}
}
\end{ex} \haidong{20}
\loai{Khoảng cách xa nhất}
\begin{ex}%[3P2G9-2][Loai1]
Trên một sợi dây đàn hồi đang có sóng dừng ổn định với khoảng cách giữa hai nút sóng liên tiếp là 24 cm. Biên độ bụng sóng là 3 cm. Gọi N là vị trí của một nút sóng; C và D là hai phần tử trên dây ở hai bên của N và có vị trí cân bằng cách N lần lượt là 8 cm và 4 cm. Khoảng cách cực đại giữa C và D trong quá trình dao động \textbf{gần nhất} với giá trị nào sau đây?
\choice
{15 cm}
{\True 12 cm}
{10 cm}
{18 cm}
\loigiai{
\begin{itemize}
\item Lấy gốc tọa độ ở nút N ta có: $A_C=A_b. \left| \sin \dfrac{2\pi x_C}{\lambda } \right|=3\left| \sin \dfrac{2\pi . 8}{24. 2} \right|=\dfrac{3\sqrt{3}}{2}\ cm;$
$a_D=A_b. \left| \sin \dfrac{2\pi x_D}{\lambda } \right|=3\left| \sin \dfrac{2\pi (-4)}{48} \right|=1. 5\ cm$.
\item Suy ra khoảng cách lớn nhất khi C, D đều ở biên theo Pitago: $\sqrt{(8+4)^2+(\dfrac{3\sqrt{3}}{2}+1. 5)^2}\approx 12,68\ cm$.
\end{itemize}
}
\end{ex} \haidong{20}
\loai{Các điểm cùng biên độ cách đều nhau}
\begin{ex}%[3P2K9-2][Loai1] 
Trên một sợi dây đàn hồi chiều dài 1,6 m, hai đầu cố định và đang có sóng dừng. Quan sát trên dây thấy có các điểm không phải bụng cách đều nhau những khoảng 20 cm luôn dao động cùng biên độ $A_0$. Số bụng sóng trên dây là
\choice
{\True 4}
{8}
{6}
{5}
\loigiai{
\begin{itemize}
\item Các điểm không phải bụng có cùng biên độ $A_0$ mà cách đều nhau một khoảng $\Delta x$ thì $A_0=\dfrac{A_{\max}}{\sqrt{2}}$ và $\Delta x=\dfrac{\lambda}{4}$.
\item Ta có: $\dfrac{\lambda}{4}=0,2\ m\Rightarrow \lambda =0,8\ m\Rightarrow sb=\dfrac{AB}{0,5\lambda}=\dfrac{1,6}{0,5.0,8}=4$.
\end{itemize}
}
\end{ex} \haidong{20}

\loai{Ba điểm liên tiếp cùng biên độ}

\begin{ex}%[3P2T9-2][Loai1] %[3P2-3.3-4] 
Do sóng dừng xảy ra trên sợi dây. Các điểm dao động với biên độ 3 cm có vị trí cân bằng cách nhau những khoảng liên tiếp là 10 cm hoặc 20 cm. Biết tốc độ truyền sóng là 15 m/s. Tốc độ dao động cực đại của bụng có thể là
\choice
{$15\pi\ cm/s$}
{$150\pi\ cm/s$}
{\True $300\pi\ cm/s$}
{$75\pi\ cm/s$}

\end{ex} \haidong{20}
\begin{ex}
Một sợi dây đàn hồi dài $30\ cm$ có hai đầu cố định. Trên dây đang có sóng dừng. Biết sóng truyền trên dây với bước sóng $20\ cm$ và biên độ dao động của điểm bụng là $2\ cm$. Số điểm trên dây mà phần tử tại đó dao động với biên độ $6\ mm$ là ?
\choice
{\True 6}
{7}
{8}
{9}
\loigiai{}
\end{ex}
\begin{ex}
Sợi dây dài $84\ cm$ có sóng dừng. Ngoài 2 đầu cố định còn 6 điểm đứng yên. $f = 40\ Hz$, biên độ bụng $6\ cm$. $M$, $N$, $P$ có cùng biên độ $3\ cm$. $N$ xa $M$ nhất cùng pha, $P$ xa $M$ nhất ngược pha. Tỉ số lớn nhất giữa $MP$ và $MN$ là
\choice 
{1,1}
{3,2}
{\True 0,9}
{9}
\loigiai{
Có 6 điểm đứng yên $\Rightarrow$ 7 bó sóng $\Rightarrow \lambda = \dfrac{2 . 84}{7} = 24\ cm$\\
$MN = \dfrac{10}{3}\lambda = 80\ cm$; $\Delta u_{MN} = 0$\\
$MP = \sqrt{(\Delta x)^2 + (\Delta u)^2} = \sqrt{74^2 + 6^2} \approx 74,2\ cm$\\
$\Rightarrow \delta_{\max} = \dfrac{MP}{MN} \approx \dfrac{74,2}{80} \approx 0,9$
}
\end{ex}
\loai{Khoảng thời gian ngắn nhất giữa hai lần li độ dao động của phần tử tại B bằng biên độ dao động của phần tử tại C}

\begin{ex}%[3P2G9-2][Loai1] 
Một sợi dây đàn hồi căng ngang, đang có sóng dừng ổn định. Trên dây, A là một điểm nút, B là một điểm bụng gần A nhất, C là điểm thuộc đoạn AB sao cho $AC=\dfrac{10}{3}$ cm, với $AB=10$ cm. Biết khoảng thời gian ngắn nhất giữa hai lần mà li độ dao động của phần tử tại B bằng biên độ dao động của phần tử tại C là 0,25 s. Tốc độ truyền sóng trên dây là
\choice
{2 m/s}
{\True 0,53 m/s}
{1 m/s}
{0,25 m/s}
\loigiai{
\immini{\begin{itemize}
\item Do $AC = \dfrac{1}{3}AB$nên biên độ dao động của C là $\dfrac{1}{2}A_B$.
\item Ta biểu diễn các vị trí li độ dao động của phần tử tại B bằng biên độ dao động của phần tử tại C tại các điểm $M_1, M_2$ trên đường tròn như hình vẽ
\item $\dfrac{T}{3} = 0,25 \Rightarrow T = 0,75\ s$.
\item Mà $\lambda = 4.AB = 40\ cm $.
\item $v = \dfrac{\lambda}{T} = \dfrac{0,4}{0,75} = 0,53\ m/s$.
\end{itemize}}{\begin{tikzpicture}[scale=1,thick,>=stealth']
\tkzInit[ymin=-3,ymax=3,xmin=-3,xmax=3]\tkzClip
\tkzDefPoints{-2.5/0/x, 2.5/0/x', 0/0/o, 0/2.5/y, 0/-2.5/y'}
\tkzDefShiftPoint[o](60:2){m}
\tkzDefPointBy[projection=onto x--x'](m) \tkzGetPoint{h1}
\tkzDefShiftPoint[o](-60:2){n}
\tkzDefPointBy[projection=onto x--x'](n) \tkzGetPoint{h2}
\draw(o)circle(2cm);
\tkzDrawSegments(y,y')
\tkzDrawSegments[dashed](m,h1 n,h2)
\tkzDrawSegments[->](x,x')
\tkzDrawSegments[blue,->](o,m o,n)
\tkzDrawPoints[size=3](o,m,n,h1,h2)
\node at (2,0)[below right]{$+A_B$};
\node at (-2,0)[below left]{$-A_B$};
\node at (h1)[below]{$\dfrac{A_B}{2}$};
\node at (x')[above]{$x$};
\node at (m)[above right]{$M_2$};
\node at (n)[below right]{$M_1$};
\end{tikzpicture}}
}
\end{ex} \haidong{20}
\loai{Trạng thái của các điểm ở 2 thời điểm khác nhau}
\begin{ex}%[3P2T9-2][Loai1] 
Trên một sợi dây đàn hồi đang có sóng dừng ổn định với khoảng cách giữa hai nút sóng liên tiếp là 6 m. Trên dây có những phần tử sóng dao động với tần số 5 Hz và biên độ lớn nhất là 3 cm. Gọi N là vị trí của một nút sóng; C và D là hai phần tử trên dây ở hai bên của N và có vị trí cân bằng cách N lần lượt là 10,5 m và 7 m. Tại thời điểm $t_1$, phần tử C có li độ 1,5 cm và đang hướng về vị trí cân bằng. Vào thời điểm $t_2=t_1 +\dfrac{3}{4}$ s, phần tử D có li độ là
\choice
{$\dfrac{-3\sqrt3}{4}$ cm}
{$\dfrac{3\sqrt3}{4}$ cm}
{$\dfrac{3\sqrt2}{4}$ cm}
{\True $\dfrac{-3\sqrt2}{4}$ cm}

\end{ex} \haidong{20}


\setcounter{ex}{0}
\PhanII
\begin{ex}
Trên sợi dây đàn hồi có chiều dài $80\ cm$, người ta tạo ra sóng dừng có hình dạng được mô tả như hình bên dưới.\\
\centerline{\begin{tikzpicture}[draw=Mapcolor,very thick,>=stealth']
\pgfplotsset{width=8cm,height=5cm,compat=1.4}
\begin{axis}[
axis lines=none,
xmin=0,xmax=4,xstep=1,ymin=-3,ymax=4.5,ystep=1,
grid=none, %Có các tùy chọn: minor|major|both|none
x grid style={dashed},
y grid style={dashed},
ytick={-1.5,-1,-0.5,0,0.5,1,1.5},
xtick={0.33,0.66,...,6},
xticklabels={},
yticklabels={},
xticklabel style={below},
xlabel style={above=1 pt},
xlabel=$x(cm)$,
ylabel style={left=2pt},
ylabel={$u$}
]
\addplot[line width=1.2pt,domain=0:4,samples=200,Mapcolor]{1.5*cos(deg(0.5*pi*x-pi/2))};
\addplot[line width=1.2pt,domain=0:4,samples=200,Mapcolor]{1.5*cos(deg(0.5*pi*x+pi/2))};
\draw[<->,very thick] (axis cs: 0, -2)--node[midway,below]{L}(axis cs: 4, -2);
\draw (axis cs: 1, 0)circle (1pt)node[below]{A};
\draw(axis cs: 3, 0)circle (1pt)node[below]{B};
\draw[<->,very thick] (axis cs: 1, 2)--++(90:0.4cm)--node[midway,fill=white]{$\dfrac{\lambda}{2}$}++(0:3cm)--++(-90:0.4cm);
\end{axis}
\end{tikzpicture}}
\choiceTF[t]
{\True Trên dây có 2 bó sóng}
{Bước sóng của sóng trên sợi dây là $60\ cm$}
{Nếu sóng truyền trên dây có tần số $50\ Hz$ thì tốc độ của sóng là $40\ cm/s$}
{\True Hai điểm A và B cách nhau $40\ cm$}
\loigiai{
\begin{itemchoice}
\itemch
\itemch
\itemch
\itemch
\end{itemchoice}}
\end{ex}
\begin{ex}
Một sợi dây $AB$ dài $100\ cm$ căng ngang, đầu $B$ cố định, đầu $A$ gắn với một nhánh của âm thoa dao động điều hòa với tần số $40\ Hz$. Trên dây $AB$ có một sóng dừng ổn định, $A$ được coi là nút sóng. Tốc độ truyền sóng trên dây là $20\ m/s$. Biên độ của bụng là $4\ cm$.
\choiceTF[t]
{Khoảng cách giữa hai bụng sóng liên tiếp bằng $50\ cm$}
{Tổng số bụng sóng và nút sóng bằng $8$}
{\True $M$, $N$ là hai phần tử dao động điều hòa có vị trí cân bằng cách đầu $A$ lần lượt $20\ cm$ và $30\ cm$. Tỉ số giữa biên độ dao động của $M$ và biên độ dao động của $N$ là $1$}
{\True Trên sợi dây, khoảng cách lớn nhất giữa hai điểm dao động cùng biên độ $2\ cm$ và ngược pha là $95,83\ cm$}
\loigiai{
\begin{itemchoice}
\itemch Sai
\itemch Sai
\itemch Đúng
\itemch Đúng
\end{itemchoice}
}
\end{ex}
\begin{ex}
Trên một sợi dây căng thẳng nằm ngang hai đầu cố định có chiều dài $\ell = 1,2\ m$ đang có sóng dừng. Cho biết mỗi điểm bụng sóng dao động theo phương thẳng đứng với biên độ $8\ mm$, khoảng cách giữa hai nút sóng liên tiếp là $0,3\ m$.
\choiceTF[t]
{\True Bước sóng trên sợi dây có giá trị là $0,6\ m$}
{Trên sợi dây có tổng số $4$ nút sóng}
{Bụng sóng là những điểm trên sợi dây không dao động}
{\True Trên sợi dây có $8$ điểm dao động với biên độ $5\ mm$}
\loigiai{
\begin{itemchoice}
\itemch
\itemch
\itemch
\itemch
\end{itemchoice}
}
\end{ex}
\begin{ex}
Trên một sợi dây đàn hồi dài $1,2\ m$ đang có sóng dừng ổn định với tần số $10\ Hz$, người ta thấy ngoài hai đầu dây cố định còn có 3 điểm khác luôn đứng yên.
\choiceTF[t]
{\True Sóng dừng trên dây là sự giao thoa của sóng tới và sóng phản xạ lan truyền trên dây, hai sóng này cùng biên độ, cùng tần số}
{\True Khoảng cách ngắn nhất giữa một cực đại và một cực tiểu bằng $15\ cm$}
{Tốc độ truyền sóng trên dây là $60\ m/s$}
{Khoảng cách giữa 3 nút sóng liên tiếp $30\ cm$}

\loigiai{
\begin{itemchoice}
\itemch Bản chất sóng dừng là sự giao thoa của sóng tới và sóng phản xạ cùng biên độ, cùng tần số.
\itemch Khoảng cách giữa cực đại và cực tiểu gần nhau nhất là $\dfrac{\lambda}{4} = \dfrac{60}{4.10} = 15\ cm$.
\itemch Với $4$ bụng thì $\lambda = \dfrac{2L}{n} = \dfrac{2.1{,}2}{4} = 0{,}6\ m$, suy ra $v = f\lambda = 10 . 0{,}6 = 6\ m/s$.
\itemch Ba nút liên tiếp cách nhau $0{,}6/2 = 0{,}3\ m = 30\ cm$, nhưng trong đề bài là "khoảng cách giữa 3 nút liên tiếp", tức là $2$ đoạn, bằng $60\ cm$.
\end{itemchoice}
}
\end{ex}

\setcounter{ex}{0}
\PhanIII
\begin{ex}
Trên một sợi dây đàn hồi đang có sóng dừng với biên độ dao động của các điểm bụng là $a$. $M$ là một phần tử dây dao động với biên độ $0,5a$. Biết vị trí cân bằng của M cách điểm nút gần nó nhất một khoảng $2\ cm$. Sóng truyền trên dây có bước sóng là bao nhiêu cm?
\shortans[oly]{24}
\loigiai{}
\end{ex}
\begin{ex}
Hai sóng hình sin cùng bước sóng $\lambda$, cùng biên độ $a$ truyền ngược chiều nhau trên một sợi dây với cùng vận tốc $20\ cm/s$ tạo ra sóng dừng. Biết hai thời điểm gần nhất mà dây duỗi thẳng là $0,5\ s$. Bước sóng $\lambda$ là bao nhiêu cm?
\shortans[oly]{20}
\end{ex}
\textit{Sử dụng dữ kiện sau cho Câu 3 và Câu 4:} Quan sát một hệ sóng dừng trên dây đàn hồi, ta thấy với M là một nút sóng và N là một bụng sóng kế cận thì khoảng cách $MN = 10\ cm$ cho biết bề rộng của một bụng sóng là 4 cm . 
\begin{ex}
Biên độ dao động của sóng là bao nhiêu cm?
\shortans[oly]{1} 
\end{ex}
\begin{ex}
Biên độ dao động của điểm I là trung điểm MN là bao nhiêu cm?
\shortans[oly]{1,4} 
\loigiai{
Vì biên độ sóng tại bụng là 2a nên bề rộng bụng sóng là 4 a .
Suy ra biên độ sóng là $a = \dfrac{4,0}{4} = 1,0$ cm
Biên độ dao động của I là trung điểm $MN: a_{I} = 2 a \left| \sin \dfrac{2 \pi d}{\lambda} \right| = 2 . 1,0 . \left| \sin \dfrac{2 \pi . \dfrac{\lambda}{8}}{\lambda} \right| = 2 \sin \dfrac{\pi}{4} = 2 . \dfrac{\sqrt{2}}{2} = 1,4$ cm
}
\end{ex}
\begin{ex}
\impicinpar{Sóng ngang có tần số f truyền trên một sợi dây đàn hồi rất dài, với tốc độ $3\ m/s$. Xét hai điểm M và N nằm trên cùng một phương truyền sóng, cách nhau một khoảng x. Đồ thị biểu diễn li độ sóng của M và N cùng theo thời gian t như hình vẽ. Biết $t_{1}=0,05\ s$. Gọi $d$ là khoảng cách giữa hai phần tử M và N tại thời điểm $t_{2}$; $\ell_{\max}$ là khoảng cách lớn nhất giữa hai phần tử M và N trong quá trình dao động. Tỉ số $\dfrac{d}{\ell_{\max}}$ là bao nhiêu (làm tròn kết quả đến chữ số hàng phần mười)?}
{\begin{tikzpicture}[draw=Mapcolor,very thick,>=stealth']
	\pgfplotsset{width=7cm,height=4.5cm,compat=1.4}
	\begin{axis}[
	axis lines=middle, xmin=0, xmax=2*pi, xstep=1,
	ymin=-4.5, ymax=5, ystep=0.1,
	grid=major,%Vẽ chế độ lưới theo trục tọa độ Oxy,
	x grid style={dashed, black},
	axis line style={very thick,Mapcolor}, % <--- dòng thêm vào
	xlabel=$t~(s)$,
	xlabel style={above},
	xtick ={1,2,3,4,5,6,7,17/3,8},
	xticklabels={},
	y grid style={dashed, black},
	extra y ticks={0},
	extra y tick label={0},
	ylabel=$u~(mm)$,
	ylabel style={above},
	ytick={-4,-2,0,2, 4},
	yticklabels={$-4$,$-2$,,2,4},]
	\addplot[line width=1.2pt, domain=0:1.9*pi, samples=400, Mapcolor]{4*cos(deg(0.5*pi*x))};
	\addlegendentry{$u_N$}
	\addplot[line width=1.2pt,dashed, domain=0:1.9*pi, samples=400, Mapcolor]{4*cos(deg(0.5*pi*x-pi/3))};
	\addlegendentry{$u_M$}
	\addplot[Mapcolor,mark=*,mark options={fill=white,mark size=1.5pt}, nodes near coords={$t_1$}]
	coordinates {(3, -4)};
	\addplot[Mapcolor,mark=*,mark options={fill=white,mark size=1.5pt}, nodes near coords={$t_2$}]
	coordinates {(17/3, -4)};
	\end{axis}
\end{tikzpicture}}
\shortans[oly]{0,9}
\loigiai{  }
\end{ex}
\begin{ex}
Một sợi dây đàn hồi căng ngang, đang có sóng dừng ổn định. Trên dây, $A$ là một điểm nút, $B$ là vị trí cân bằng của một điểm bụng gần $A$ nhất, $C$ là trung điểm của $A B$, với $A B=10\ cm$. Biết khoảng thời gian ngắn nhất giữa hai lần mà li độ dao động của phần tử tại $B$ bằng biên độ dao động của phần tử tại $C$ là $0,2\ s$. Tốc độ truyền sóng trên dây là bao nhiêu $m/s$ (làm tròn kết quả đến chữ số hàng phần mười)?
\shortans[oly]{0,5}
\loigiai{
\item Khoảng cách giữa một nút sóng và bụng sóng gần nhất là $\dfrac{\lambda}{4}$. 
\item Vậy $AB = \dfrac{\lambda}{4} = 10 cm \Rightarrow \lambda = 40 cm = 0,4\ m$.
\item Biên độ dao động tại bụng sóng B là $A_B = A_{max}$.
\item Điểm C là trung điểm của AB, cách nút A một khoảng là $AC = \dfrac{AB}{2} = \dfrac{10 cm}{2} = 5\ cm$.
\item Biên độ dao động tại điểm cách nút một khoảng $d$ là $A = A_{max} |\sin\left(\dfrac{2\pi d}{\lambda}\right)|$.
\item Vậy biên độ dao động tại C là $A_C = A_{max} |\sin\left(\dfrac{2\pi . 5 cm}{40 cm}\right)| = A_{max} |\sin\left(\dfrac{\pi}{4}\right)| = A_{max} \dfrac{\sqrt{2}}{2}$.
\item Li độ dao động của phần tử tại B là $x_B = A_{max} \cos(\omega t)$.
\item Li độ dao động của phần tử tại C là $x_C = A_C \cos(\omega t + \varphi_C) = A_{max} \dfrac{\sqrt{2}}{2} \cos(\omega t + \varphi_C)$. \item Do C gần nút hơn bụng nên dao động cùng pha với điểm gần nút, và ngược pha với điểm gần bụng. 
\item Vậy $\varphi_C = 0$.
$x_C = A_{max} \dfrac{\sqrt{2}}{2} \cos(\omega t)$.
\item Theo đề bài, khoảng thời gian ngắn nhất giữa hai lần mà $|x_B| = |A_C|$ là $0,2 s$.
\item $|A_{max} \cos(\omega t)| = |A_{max} \dfrac{\sqrt{2}}{2}| \Leftrightarrow |\cos(\omega t)| = \dfrac{\sqrt{2}}{2}$.
\item Phương trình này có các nghiệm $\omega t = \pm \dfrac{\pi}{4} + k\pi$. 
\item Khoảng thời gian ngắn nhất giữa hai nghiệm liên tiếp là $\Delta t = \dfrac{\pi/2}{\omega} = \dfrac{T}{4} = 0,2 s$.
\item Vậy chu kì dao động $T = 4 . 0,2 s = 0,8 s$.
\item Tốc độ truyền sóng trên dây là $v = \dfrac{\lambda}{T} = \dfrac{0,4 m}{0,8 s} = 0,5 m/s$.
}
\end{ex}






